\sgaStableTarget{sga6.e13.rem.5-3}\textbf{Remark 5.3.} Let @@M0@@ be an algebraically closed field, and @@M1@@ a proper smooth surface over @@M2@@. Then the second Betti number @@M3@@ can be interpreted as @@M4@@, where @@M5@@, and @@M6@@ = the second Chern class = the Euler-Poincar\'e characteristic = the intersection number of the diagonal of @@M7@@ with itself; thus @@M8@@ has an interpretation independent of cohomology. The inequality @@M9@@ in this form for @@M10@@ is due (before \'etale cohomology) to IGUSA, who used deep results on the structure of the tame fundamental group of an algebraic curve minus a few points.

\sgaStableTarget{sga6.e13.thm.5-4}\textbf{Theorem 5.4.} \emph{Let @@M11@@ be an algebraically closed field, and @@M12@@ a proper @@M13@@-scheme. Then there exist finitely many closed integral curves @@M14@@ in @@M15@@ with the following property: a subset @@M16@@ of the scheme @@M17@@ is quasi-compact if and only if the intersection numbers @@M18@@ for @@M19@@ remain bounded. Moreover, we may take for @@M20@@ the rank of the N\'eron-Severi group @@M21@@.}

Indeed, since @@M22@@ is an open quasi-compact subgroup of @@M23@@ (4.7), @@M24@@ is quasi-compact if and only if the corresponding set of points in the vector space @@M25@@ is finite. But @@M26@@-equivalence is equal to numerical equivalence (4.6 (i) @@M27@@ (iii)), so we may choose among the curves in @@M28@@ curves that map to a basis of the dual space; hence the assertion.

